\documentclass[11pt]{article}
\usepackage[margin=1in]{geometry}
\usepackage{graphicx}
\usepackage{amsmath,amssymb}
\usepackage{xcolor}
\usepackage[hidelinks]{hyperref}
\usepackage{parskip}

\title{\textbf{Coupling discrete Turing reaction--diffusion to a 3D vertex model:\\
what transfers cleanly, and what needs force balance}}
\author{Notes from a \texttt{plexus2} prototype of Okuda et al.\ (2018)}
\date{July 2026}

\begin{document}
\maketitle

\begin{abstract}
\noindent
We report the lessons from building a \texttt{plexus2} prototype that couples a \emph{discrete}
Turing reaction--diffusion (RD) system to a 3D vertex / self-propelled-Voronoi model on a growing
epithelial vesicle, following Okuda \emph{et al.}, \emph{Sci.\ Rep.} \textbf{8}:2386 (2018). Four
findings stand out. \textbf{(1)} When growth is morphogen-\emph{independent}, the RD is a passive
\emph{colouring} of the cells: it never touches the mechanics. Apparent ``noise'' in Gray--Scott
patterns was therefore not a chemistry problem but a \emph{timestep} artefact --- the coarse step
Gray--Scott needs to develop ($\Delta t\!=\!0.35$) destabilises the vertex mechanics, which had to
be detuned to stay stable. Running the mechanics at a fine, stable step ($\Delta t\!=\!0.02$) while
\emph{time-rescaling} the RD by a single rate multiplier recovers a clean epithelium carrying the
coral / labyrinth / holes patterns (cell-shape roughness $0.04$ vs.\ $0.24$). This is exactly the
chemistry-vs-mechanics timescale separation ($\Delta t_m \ll \Delta t_v$) of the original model.
\textbf{(2)} A latent operator bug --- a missing \texttt{SphericalVoronoi} import silently swallowed
by a broad \texttt{except} --- had made the tangential (Lloyd) relaxation a \emph{total no-op} for
the operator's entire life; the cell uniformity seen in earlier renders came \emph{entirely} from a
render-side cosmetic, which masked the dead physics. \textbf{(3)} Reproducing morphogen-driven
\emph{tubulation} (their Fig.~5) is qualitatively harder: an overdamped, single-step vertex
relaxation \emph{without} force-balance iteration or topological reconnection converts localised
growth into \emph{global} chaotic buckling, uncorrelated with the morphogen
($\mathrm{corr}(a,r)\!\approx\!0$), rather than local tubes. \textbf{(4)} A working alternative
sidesteps the fragile shell mechanics entirely: amplifying the (already present, but weak)
activator-dependent \emph{division} into a sharp $\sim\!4{:}1$ ON/OFF contrast on the fast, stable
spherical mechanics produces \emph{localised} extrusion that tracks the morphogen
($\mathrm{corr}(a,r)\!=\!0.31$ vs.\ $\approx\!0$ for the shell), while the monolayer stays clean --
the high-activator domains bulge out into lobes.
\end{abstract}

\section{The model in one paragraph}
One \texttt{cell} set carries three independently integrated state blocks: position \texttt{pos}
(the vertex coordinate), morphogens \texttt{chem}\,$=\,[a,h]$ (a Gray--Scott or Brusselator RD on
the cell--cell adjacency graph, so mechanical deformation reshapes the diffusion), and a target
size \texttt{v0} that a Hill-saturating growth law inflates. Division wakes a dormant buffer slot
when \texttt{v0} reaches its threshold. Mechanics is an overdamped shape energy
$U=\sum_j K_V(V_j-V_j^{0})^2 + K_S S_j$ over ghost-bounded Voronoi cells; the force is
$-\nabla U$ by autodiff through the tetrahedral circumcentres. Each of these is a \texttt{plexus2}
operator (a typed morphism on the set) composed by a schedule; different operators run at different
cadences, which is how the multiple timescales are expressed.

\section{What transferred cleanly: Fig.~4 (patterns on a growing vesicle)}
In Fig.~4 growth is \emph{activator-independent} (homogeneous), so the RD only colours the cells.
The naive Gray--Scott vesicle looked badly corrupted --- a lumpy, non-epithelial surface --- and it
was tempting to blame the reaction--diffusion. It was not: the mechanics was simply being integrated
at Gray--Scott's coarse $\Delta t\!=\!0.35$, far outside the vertex model's stability window. The fix
is to \emph{decouple the timescales}: keep the mechanics at the fine step that makes it clean
($\Delta t\!=\!0.02$, the same parameters that produce a nice Brusselator vesicle) and speed the RD
up by a rate multiplier $\rho$ (with the diffusion coefficient scaled to match), so the identical
pattern develops in the same number of frames. Roughly $\rho\!\approx\!\Delta t_{\text{coarse}}/
\Delta t_{\text{fine}}$; the only ceiling is the explicit-diffusion CFL limit
($\Delta t\,d_h\,\chi\,\bar{k}\lesssim 0.4$), which the delicate low-feed labyrinth regime hits first
and which must be respected per regime. The result (Fig.~\ref{fig:coral}) is a uniform epithelium
carrying the branching coral pattern, with shell roughness $0.04$ against $0.24$ for the coarse-step
version. The take-away is general: \emph{when a signal does not feed back on shape, do not let its
integration timescale dictate the mechanics'}.

\begin{figure}[h]
\centering
\includegraphics[width=0.92\textwidth]{archive/fig4_coral_v/strip.png}
\caption{\textbf{Gray--Scott coral on the vesicle mechanics} (\texttt{fig4\_coral\_v}), at
$0/50/100\%$ of the run. Top: closed external view; bottom: cross-section through the monolayer and
lumen. Cells are uniform (near-hexagonal) and the shell is smooth; the reaction--diffusion is a pure
colouring of a clean epithelium. $N\!:\!600\!\to\!3000$, roughness $0.04$.}
\label{fig:coral}
\end{figure}

\section{A cautionary tale: the silent no-op}
The uniform cells above depend on a tangential (Lloyd) relaxation that moves each cell toward its
spherical-Voronoi centroid. That operator had \emph{never once} executed: the module imported
\texttt{Delaunay, ConvexHull, cKDTree} but not \texttt{SphericalVoronoi}, so every call raised
\texttt{NameError}, which fell into a blanket \texttt{except Exception: return zeros}. The knob was
inert --- \texttt{k\_lloyd} $0.3/1.0/2.0$ gave byte-identical output. The uniformity previously
visible came entirely from a \emph{render-side} Lloyd pass (a cosmetic applied only when drawing).
The moment we honoured ``uniformity should come from physics, not the render'' and drew the true
cells, the dead operator was exposed. Two lessons: a broad \texttt{except} around a geometry call can
hide a programming error for the lifetime of a component; and a cosmetic post-process can mask a
missing mechanism indefinitely. The fix was one import plus narrowing the handler to the genuine
Qhull degeneracies; the operator was also made timestep-independent (move a fixed \emph{fraction}
toward the centroid per tick, capped at $0.7\times$ the cell spacing).

\section{What did not transfer: Fig.~5 (tubulation)}
Fig.~5 makes the activator a \emph{mitogen}: high-activator cells grow, and on a closed sheet the
locally excess area must buckle out of plane into tubes --- no explicit extrusion force. We
implemented the coupling (a Hill switch from activator concentration to target-volume growth) on the
apical/basal ``shell'' mechanics, and swept $\sim\!15$ configurations of domain size, growth rate,
bending, surface tension, and lumen pressure. \emph{None} produced clean tubes. The behaviour is
bimodal: gentle growth or stiff bending gives a perfectly smooth sphere (no buckling), while enough
growth gives \emph{global} chaotic buckling that tears the monolayer into a blob. Decisively, the
correlation between activator concentration and radial protrusion is $\approx\!0$ in every case: the
bulges are \emph{not} where the activator is --- it is generic Euler buckling of an over-grown,
lumen-pressurised shell, not localised extrusion.

Three mechanics ingredients are missing, all present in the original: (i) \emph{force-balance
iteration} --- Okuda solves the overdamped balance to residual $<\!10^{-5}$ each $\Delta t_v$,
whereas a single explicit step lets individual cells spike rather than forming coherent folds;
(ii) \emph{topological reconnection} (RNR / T1), without which cells jam instead of flowing into a
tube; and (iii) a tractable per-cell tessellation --- our \texttt{ConvexHull}-per-cell shell becomes
pathologically slow, even hanging, exactly in the heavily-deformed regime tubes require. Clean
tubulation is thus gated on a vectorised, force-balance-iterated shell mechanics plus a reconnection
operator, not on more parameter tuning.

\section{Field notes: what we learned about the model}
Grouped by subsystem. The numbers are the values that actually worked --- or failed --- in this
prototype; they are the hard-won part and the main reason for this document.

\paragraph{Timescales --- the master knob.}
\begin{itemize}\itemsep2pt
\item Gray--Scott needs a large step ($\Delta t\!\approx\!0.35$) to develop in $\sim\!10^4$ frames;
the overdamped vertex model is stable only for $\Delta t\!\lesssim\!0.05$. Irreconcilable at a
single step.
\item Resolution: run mechanics at $\Delta t\!=\!0.02$ and multiply the RD rate by
$\rho\!\approx\!\Delta t_{\text{coarse}}/\Delta t_{\text{fine}}\!\approx\!17.5$, scaling the
diffusion prefactor $\chi$ by the \emph{same} $\rho$ so the reaction:diffusion ratio --- and thus
the pattern --- is preserved. A pure time-rescaling: identical pattern, stable mechanics.
\item Explicit graph-diffusion is stable only while
$\Delta t\,d_h\,\chi\,\langle k\rangle\!\lesssim\!0.4$. With $d_h\!=\!0.16,\ \langle k\rangle\!
\approx\!14,\ \Delta t\!=\!0.02$ this caps $\chi\!\lesssim\!9$. Robust high-feed coral
($F\!=\!0.058$) tolerates $\chi\!=\!9.8,\rho\!=\!35$; the delicate low-feed labyrinth
($F\!=\!0.029$) \emph{diverges to NaN} there and needs $\chi\!=\!5.6,\rho\!=\!20$ over more frames.
Tune the rescaling \emph{per regime}.
\end{itemize}

\paragraph{Tangential (Lloyd) relaxation.}
\begin{itemize}\itemsep2pt
\item Make it timestep-independent: move a fixed fraction \texttt{k\_lloyd} of the way to the
spherical-Voronoi centroid per tick, then divide by $\Delta t$ to undo the engine's
$x\!\mathrel{+}=\!\Delta t\,v$. \texttt{k\_lloyd} is then a clean fraction ($1.0$ = one full Lloyd
step/tick; $1.5$ slightly overshoots and converges faster). Cap the per-tick move at
$0.7\times$ mean spacing so a defect cell cannot jump neighbours and tangle the mesh.
\item It uses \texttt{SphericalVoronoi} of the cell \emph{direction} vectors, which degenerates once
the surface is non-spherical --- so it is valid on the vesicle, \emph{not} on Fig-5 tubes.
\end{itemize}

\paragraph{Morphogen $\to$ growth (the coupling).}
\begin{itemize}\itemsep2pt
\item $\lambda=\lambda_{\text{ref}}\big(\rho + g\,\mathrm{Hill}(a)\big),\ \mathrm{Hill}(a)=
a^{\alpha}/(a^{\alpha}+a_{\text{sw}}^{\alpha})$. $\rho$ is the OFF multiple, $\rho+g$ the ON one.
\item $a_{\text{sw}}$ \emph{must} sit inside the activator's actual range or the switch never fires.
The Gray--Scott activator peaks near $0.4$, so $a_{\text{sw}}\!=\!1.0$ gives
$\mathrm{Hill}\!\approx\!0.16$ (a feeble $16\%$ contrast); $a_{\text{sw}}\!\approx\!0.2$ makes it
bite. $\alpha\!=\!4$ gives a near-step switch.
\item Even ``homogeneous'' Fig-4 growth ($\rho\!=\!1$) carries the Hill term, so high-activator
domains already bulge weakly. Amplify with $g$ ($\rho\!=\!0.5,g\!=\!1.5\!\Rightarrow\!4{:}1$) for
real localised proliferation.
\end{itemize}

\paragraph{Division.}
\begin{itemize}\itemsep2pt
\item \texttt{divide\_2x} fires when $v^0$ reaches $\texttt{ratio}\!\times\!$base
($v_{\text{th}}\!=\!\tfrac43 v_{\text{ref}}$); the daughter wakes a dormant buffer slot, both reset
to base, and the daughter inherits the mother's \texttt{chem}. Faster growth $\Rightarrow$ earlier
division $\Rightarrow$ local density $\Rightarrow$ local area excess --- the buckling driver.
\item \texttt{reset\_noise} spreads daughter target sizes so the cycle stays asynchronous (no
division waves). Fixed-buffer occupancy: division stops when the buffer is full.
\end{itemize}

\paragraph{Mechanics --- two flavours; know which.}
\begin{itemize}\itemsep2pt
\item \emph{Spherical-ghost} \texttt{voronoi\_tension\_3d}: fast (vectorised, $\sim\!32\times$ over
the per-cell loop), stable, keeps a clean monolayer --- but the global ghost sphere \emph{fights}
local deformation, so at most mild bumps. Right for Fig 4 and the differential-growth route.
\item \emph{Per-cell apical/basal shell} \texttt{voronoi\_tension\_shell}: allows local buckling but
needs a scipy \texttt{ConvexHull} per cell; under heavy deformation this is pathologically slow and
can \emph{hang} (observed 77-min hang, 84-min no-finish). Not viable for tubes un-vectorised.
\item Stiff volume vs.\ surface, $K_V/K_S\!\approx\!10$ (incompressible cells). Bending
$k_{\text{bend}}$ HIGH ($0.6$--$0.8$) forbids out-of-plane curvature $\to$ no buckle;
$k_{\text{bend}}\!\approx\!0.45$ permits it.
\end{itemize}

\paragraph{Fig-5 buckling phenomenology (all negative, all informative).}
\begin{itemize}\itemsep2pt
\item Bimodal: gentle growth / stiff bending $\to$ smooth sphere; enough growth $\to$ global chaotic
buckle that tears the monolayer. No clean-tube window in between.
\item \texttt{lumen\_pressure} holds the closed vesicle rigid, so area is stored until it releases as
one global Euler buckle at a random site --- morphology decoupled from the pattern. Removing the
lumen did \emph{not} localise it ($\mathrm{corr}(a,r)$ stayed $\approx\!0$).
\end{itemize}

\paragraph{Diagnostics that earned their keep.}
\begin{itemize}\itemsep2pt
\item \emph{roughness} $=\mathrm{std}/\mathrm{mean}$ of cell radii ($0.04$ clean, $0.24$ corrupted);
\emph{\% hexagons} (closed sphere tops out $\sim\!70\%$; a bounded 2D disc lower, its boundary ring
never hexagonal); \emph{PROTR} $=$ 99th-percentile$/$median radius ($\approx\!1$ sphere, $>\!1.4$
protrusions); \emph{$\mathrm{corr}(a,r)$} $=$ the decisive extrusion-is-localised test.
\end{itemize}

\paragraph{Performance / engineering.}
\begin{itemize}\itemsep2pt
\item Multi-rate composition via an engine \texttt{every:K} gate plus an intra-operator topology
cache (\texttt{retess\_every}): mechanics reuses the tessellation for $K$ ticks, pattern preserved,
$\sim\!1.7\times$ at $K\!=\!10$.
\item \texttt{torch.compile} gives $\sim\!50\times$ on the fixed-$N$ energy but recompiles when $N$
grows (division) --- keep it off for growing runs. CPU/scipy is the real bottleneck; on the cluster,
one dedicated node per generation ($\sim\!3\times$ over contended local runs).
\end{itemize}

\section{Discussion: a cleaner route to extrusion}
The negative result above is informative. Localised growth on a \emph{closed, pressurised} vesicle
tends to store area until it releases as a single global buckle at a random site --- morphology
decoupled from the pattern. A better substrate for \emph{local} extrusion is the fast, stable
\emph{spherical} mechanics used for Fig.~4 (it never breaks the monolayer), driven by \emph{localised
proliferation} rather than global inflation.

Strikingly, that mechanism is already latent in the Fig.~4 runs. Even ``homogeneous'' growth carries
a Hill term, $\dot v^0 = \lambda_{\text{ref}}\big(\rho + \mathrm{Hill}(a)\big)$, so high-activator
cells already divide somewhat faster and their domains bulge slightly outward --- the small
extrusions one can see in the coral cross-section. They are weak only because the switch was mis-set
($a_{\text{sw}}\!=\!1.0$ while the Gray--Scott activator peaks near $0.4$, so $\mathrm{Hill}\approx
0.16$). Introducing a gain, $\dot v^0 = \lambda_{\text{ref}}\big(\rho + g\,\mathrm{Hill}(a)\big)$
with $\rho\!=\!0.5,\,g\!=\!1.5$ and $a_{\text{sw}}$ moved into the activator's actual range, turns
this into a sharp $2\times$-ON / $0.5\times$-OFF ($4{:}1$) \emph{division} contrast that concentrates
proliferation --- and hence out-of-plane buckling --- at the morphogen domains, on mechanics that
stay stable and cheap.

This works (Fig.~\ref{fig:ext}). The high-activator domains bulge outward into coherent lobes, the
monolayer stays intact (roughness $0.058$, cells $54\%$ hexagonal), and crucially the deformation now
\emph{tracks the signal}: $\mathrm{corr}(\text{activator},\text{radius})\!=\!0.31$, against
$\approx\!0$ for every shell-mechanics attempt. It is not yet the fully elongated tubes of the
original --- for that the force-balance and reconnection ingredients of \S5 are still needed --- but
it is genuine, localised, morphogen-driven extrusion on a substrate that never breaks. This reframes
tubulation as \emph{patterned proliferation} on a robust substrate rather than differential growth on
a fragile deformable shell.

\begin{figure}[h]
\centering
\includegraphics[width=0.92\textwidth]{archive/fig4_coral_ext/strip.png}
\caption{\textbf{Locally-favoured extrusion} (\texttt{fig4\_coral\_ext}): the same coral RD on the
fast spherical mechanics, but with a $4{:}1$ activator-driven division contrast
($\rho\!=\!0.5,g\!=\!1.5,a_{\text{sw}}\!=\!0.2$). The high-activator (dark) domains extrude outward
into lobes while low-activator regions stay flat; the cross-section (bottom) shows the monolayer
lobing out where the signal is high, and remaining a clean single sheet.
$\mathrm{corr}(a,r)\!=\!0.31$.}
\label{fig:ext}
\end{figure}

\paragraph{Can the lobes be elongated? (a cluster sweep).}
We ran a parameter sweep on the L4 cluster to push the lobes toward tubes. Two findings. \emph{(i)
``Run longer'' is the wrong lever, and it is dangerous.} Letting the tissue over-grow (bigger buffer,
more frames) both \emph{relaxes} the protrusions back to a sphere once growth saturates
($\mathrm{protr}_{\max}$ falls $1.27\!\to\!1.08$, $\mathrm{corr}\to0$) and, worse, packs cells so
densely at the activator domains that the explicit-diffusion \emph{CFL is violated} --- the RD blows
up to NaN (a silent failure we now catch with an \texttt{act\_nan} flag). The extrusion lives in a
finite $N\!\sim\!3000$ window. \emph{(ii) Sharper division contrast elongates --- but toward
spikes.} At the productive scale, raising the ON/OFF contrast from $4{:}1$ to $\sim\!6.7{:}1$
($\rho\!=\!0.3,g\!=\!1.7$) is the only lever that clearly extends the tips
($\mathrm{protr}_{\max}\!=\!1.46$, the tallest we reached), and the tips stay red (activator-led). But
Fig.~\ref{fig:contrast} shows what ``taller'' means here: the broad lobes sharpen into pointed
\emph{fingers} --- individual cells and small clusters shooting out --- rather than coherent tubes.
Lowering surface tension ($K_S\!=\!0.4$--$0.25$) helps mildly ($\mathrm{protr}_{\max}\!\approx\!1.2$--$1.27$).
Nor does bending rescue it: adding bending smooths the spikes but \emph{flattens the whole extrusion}
--- both $\mathrm{protr}_{\max}$ and $\mathrm{corr}$ fall ($1.44\!\to\!1.36$, $0.23\!\to\!0.07$). Across
the whole sweep the elongation \emph{plateaus at $\mathrm{protr}_{\max}\!\approx\!1.44$}: you cannot get
tall, clean, and localised at once on this mechanics --- push contrast and it spikes, add bending and it
flattens. This is the same wall as \S5 from the other side: the fast spherical SPV mechanics can
\emph{localise} growth beautifully (best balanced config \texttt{ext3\_contrast}: $\mathrm{protr}_{\max}
\!=\!1.44$, $\mathrm{corr}\!=\!0.23$, roughness $0.039$) but cannot hold a coherent extended neck. Clean
elongation is gated on the coherent vertex mechanics (shared vertices + force balance + T1), i.e.\ a true
vertex model such as Tyssue --- the natural next substrate.

\paragraph{Dividing directly on activation.} The cleanest coupling of all is to divide cells
\emph{directly} where the activator is on, rather than routing through growth-then-size (a one-line
\texttt{divide\_2x} option). The subtlety is that the per-tick division probability must be \emph{graded}
in the activator, $p=\texttt{act\_rate}\cdot a^{n}/(a^{n}+a_{\text{thr}}^{n})$: a hard threshold lets pale
cells that merely exceed it divide at the \emph{full} rate, so low-activator cells proliferate too (a flaw
we caught by eye and fixed). With the graded, sharply-gated rate (\texttt{ext5\_sharp}, $a_{\text{thr}}
\!=\!0.35,\ n\!=\!6$), proliferation is confined to the deep-activator cells --- Fig.~\ref{fig:actdiv}
shows the signature directly in the cell \emph{sizes}: the activator (dark) bands fill with small, densely
packed daughters and bulge outward, while the low-activator regions keep their large original cells and
stay flat. This is the most \emph{pattern-faithful} extrusion of the whole study, $\mathrm{corr}(a,r)
\!=\!0.42$ --- well above the best growth-driven config ($0.23$) and the broad-lobe \texttt{coral\_ext}
($0.31$) --- on a clean monolayer. The stable window is narrow (combining direct division \emph{with}
growth, or too high a rate, over-proliferates and blows the RD CFL to NaN). So direct, activation-graded
proliferation is the tightest pattern$\to$morphology coupling available on this mechanics; the residual
neck-coherence wall still caps the \emph{height} ($\mathrm{protr}_{\max}\!\approx\!1.4$), which is what the
vertex model would relieve.

\begin{figure}[h]
\centering
\includegraphics[width=0.92\textwidth]{archive/ext5_sharp/strip.png}
\caption{\textbf{Divide-on-activation, activator-graded} (\texttt{ext5\_sharp}): division probability
$\propto\mathrm{Hill}(\text{activator})$, no growth. Proliferation is confined to the activator (dark)
bands --- read it off the cell sizes: dark bands are small, dense divided cells that bulge out; pale
regions keep their large, undivided cells and stay flat. The most pattern-faithful extrusion of the study,
$\mathrm{corr}(a,r)\!=\!0.42$, on a clean single-sheet monolayer (cross-section, bottom).}
\label{fig:actdiv}
\end{figure}

\begin{figure}[h]
\centering
\includegraphics[width=0.92\textwidth]{archive/ext2_contrast/strip.png}
\caption{\textbf{Elongation by sharper contrast} (\texttt{ext2\_contrast}, $6.7{:}1$): the tallest
protrusions we reached ($\mathrm{protr}_{\max}\!=\!1.46$), tips still activator-led (red). But the
extrusions sharpen into pointed \emph{fingers}/spikes rather than coherent tubes --- the signature of
driving growth harder on a mechanics that lacks shared vertices and force balance.}
\label{fig:contrast}
\end{figure}

Two methodological points recur. First, \emph{signalling is geometry-coupled but not
mechanics-coupled unless you say so}: keeping that boundary explicit turned an intractable ``noisy
patterning'' problem into a one-line timescale fix. Second, in a typed-operator language the right
unit of change is often a new \emph{implementation} or \emph{parameter} of an existing contract (the
$4{:}1$ growth is still the \texttt{growth} operator), not a new operator --- ``prototype freely,
promote mechanisms conservatively.''

\section{The gap to Figure~5, and the Tyssue bridge}
This prototype set out to reach Okuda's tubulation (Fig.~5) and branching (Fig.~6) and did not --- but
the failure is now precisely located, and a parallel prototype (\texttt{prototype/Tyssue}) is already
most of the way across the gap.

\paragraph{The gap, in the language of the model ladder.} Confluent-tissue mechanics forms a ladder of
increasingly permissive representations, distinguished by \emph{what the degrees of freedom are}. At the
bottom sits the \emph{Self-Propelled-Voronoi} (SPV) model we used: the only DOF are the cell
\emph{centres}, and each cell's shape is not stored but \emph{read off}, every step, as the Voronoi
region of the centres (Bi 2016; Sussman 2017). One rung up is the \emph{vertex model}, whose DOF are the
shared tri-cellular \emph{junctions}: cells are generic polygons with explicit edges and line tensions,
free to take shapes no Voronoi partition admits. At the top is the full 3D vertex / deformable-cell
model (Okuda 2018; SimuCell3D), where each cell is a polyhedron with an apical (lumen-facing) and basal
surface, an enclosed volume, and a dynamically reconnecting network.

A tube is exactly what the bottom rung cannot hold. It is (i) a high-aspect-ratio junctional network far
from any Voronoi tessellation of the centres --- which the SPV forces the geometry back onto at every
step; (ii) an intrinsically 3D object with two principal curvatures and an apical/basal axis --- which a
planar point-tessellation has no coordinates for; and (iii) defined by an enclosed \emph{lumen} --- of
which the SPV has no notion. This is why our extrusion \emph{localised} superbly (graded
divide-on-activation, $\mathrm{corr}\!=\!0.42$) but its \emph{height} plateaued at $\mathrm{protr}_{\max}
\!\approx\!1.4$: with the shape slaved to the centres and no way to store a coherent neck, driving growth
harder gave single-cell spikes, not tubes. \emph{The wall is representational, not parametric.}

\paragraph{What Okuda's tube actually needs} --- four coupled ingredients, none native to the SPV.
\emph{(1) Quasi-static force balance:} the overdamped vertex equations solved to residual $<\!10^{-5}$
\emph{between} growth steps (Eq.~1; $\tau_{\text{cycle}}\!=\!50\!\gg\!\Delta t_v\!=\!0.005$).
\emph{(2) Incompressible cells $+$ surface tension:} $U=\sum_j\tfrac12 k_v(v_j-v_{\text{eq},j})^2+\kappa_s
s_j$, $k_v\!=\!10\!\gg\!\kappa_s\!=\!0.2$ (Eq.~3) --- growth of a cell's target volume must be absorbed
by shape change, not compression. \emph{(3) Reversible network reconnection (the 3D T1):} any edge below
$\Delta l_{\text{th}}\!=\!0.05$ reconnects, attempted every $\Delta t_r\!=\!1.0$ --- the topology change
that lets a bump \emph{flow} into a tube. \emph{(4) Activator$\to$growth:} a discrete Gierer--Meinhardt
mitogen (Hill $\alpha\!=\!10$ step at $\rho_{\text{sw}}\!=\!0.5$, Eqs.~7--8), the spot self-maintaining at
the growing tip. Tube diameter tracks the steady domain size $\propto\chi^{1/4}\phi^{1/2}$ (thick:
$\chi\!=\!0.1,\gamma\!=\!1$; thin: $\chi\!=\!0.01,\gamma\!=\!100$).

\paragraph{Tubes, branches, undulations are one knob apart.} The \emph{same} growing vesicle gives
undulation, tubulation, or branching depending only on the ratio of patterning rate $\gamma$ to
deformation rate $1/\tau_{\text{cycle}}$ --- the \emph{hysteresis} of the Turing pattern on the growing
domain. $\gamma\!=\!100$ (pattern outruns growth) $\to$ domains wander $\to$ \emph{undulation} (Fig.~7);
$\gamma\!=\!1$ (domains held in place) $\to$ persistent tip zones $\to$ \emph{tubes} (Fig.~5);
$\gamma\!=\!0.01$ with small domains ($\chi\!=\!0.01$) $\to$ division-induced morphogen dilution
repatterns a spot into separated sub-spots $\to$ \emph{branching} (Fig.~6). The last is a genuinely
\emph{discrete-cell} effect: because the RD is homogeneous \emph{per cell}, growth/division inject a
cell-scale discontinuity that, under high hysteresis, seeds the branch --- precisely where a discrete
vertex model does work a continuum cannot.

\paragraph{Rigidity is the precondition.} Whether a tissue \emph{can} tubulate is set by the shape-index
transition (Bi 2015): at $p_0\!=\!P_0/\sqrt{A_0}\!\approx\!3.81$ the T1 barrier vanishes and the sheet
fluidises. Below it the sheet jams (growth builds stress but the wall cannot flow); above it, it flows.
Okuda does not tune $p_0$ --- he sets the reconnection \emph{cadence} ($\Delta t_r\!\ll\!\tau_{\text{cycle}}$)
so the tissue is quasi-statically fluid. Two ratios thus select morphology: reconnection-vs-growth (must
be fluid to tube at all) and patterning-vs-growth (undulate/tube/branch). Our SPV had \emph{neither} knob
--- no explicit T1, no shape-index handle --- the same gap seen from the rigidity side.

\paragraph{The Tyssue bridge (parallel prototype).} We built \texttt{prototype/Tyssue}, a true vertex
model (AVM) as \texttt{plexus2} operators, expressly to supply the two ingredients the Voronoi route
lacks --- force-balance iteration and explicit T1 --- under the \emph{same} shape-energy contract. It is
already well across the gap. In 2D it reproduces the $p_0^*\!\approx\!3.81$ rigidity transition from
genuine force balance (residual energy collapses $\sim\!4$ orders across $p_0^*$, no timestep tuning ---
the SPV's ``master-knob'' problem gone), with validated T1, oriented division, apoptosis, growth and RD
operators. In 3D it already carries a closed epithelial \emph{vesicle} half-edge mesh
(\texttt{seed\_mesh\_3d}), a 3D shape energy with a \emph{per-cell lumen} volume (\texttt{shape\_energy\_3d}
--- the incompressibility $+$ surface tension of Eq.~3), 3D cell \emph{division} (\texttt{divide\_3d}),
quasi-static growth, and a \texttt{vesicle\_fig4} preset reproducing our Fig.~4 growth scale on the true
mesh.

\paragraph{The continuing steps to Fig.~5, then Fig.~6.} Against Okuda's ingredient list, Tyssue needs,
in order:
\begin{enumerate}\itemsep1pt
\item \textbf{A 3D reconnection --- the one genuinely missing mechanics op.} A 3D T1 / IH--HI reconnection
on the vesicle mesh (only the 2D T1 exists today), triggered below an edge-length threshold on a slow
cadence ($\Delta t_r\!\ll\!\tau_{\text{cycle}}$). This is what lets a bump \emph{extend} into a tube rather
than jam; Barton's continuous shrink-to-point edge flip is the robust recipe if a discrete threshold
proves brittle under large deformation.
\item \textbf{Couple the RD to the 3D mesh.} Tyssue's RD runs in 2D today and the vesicle grows
\emph{uniformly}; run the discrete Gierer--Meinhardt (or Gray--Scott) on the vesicle's cell adjacency,
flux $\propto$ (shared face area)/(cell volume) so pattern and geometry co-evolve (Eq.~2), reduced to the
two knobs $\chi$ (diameter) and $\gamma$ (hysteresis).
\item \textbf{Port the winning growth coupling onto the mesh.} Carry over exactly what won here ---
activator-\emph{graded} proliferation ($\dot v_{\text{eq}}\!\propto\!\mathrm{Hill}(a)$, or divide-on-activation
at the graded rate, $\mathrm{corr}\!=\!0.42$) --- with oriented (Hertwig) division. On the true mesh,
force balance and T1 should convert the SPV's spiky fingers into coherent constant-diameter tubes.
\emph{With these three, the prototype should tubulate.}
\item \textbf{Sweep $(\chi,\gamma)$.} Reproduce the morphology diagram (Fig.~8) on a $\sim\!2000$-cell
vesicle: $\gamma\!=\!100\to$ undulation, $\gamma\!=\!1\to$ tube, $\gamma\!=\!0.01,\chi\!=\!0.01\to$
branching; verify the $\chi^{1/4}$ diameter scaling.
\item \textbf{Preserve per-cell discreteness for branching.} Keep the RD field homogeneous \emph{per cell}
and run high hysteresis, so division-induced dilution injects the cell-scale perturbation that splits a
domain into a branch (Fig.~6).
\end{enumerate}
The lumen is already in \texttt{shape\_energy\_3d}; membrane bending can be added later to set thin-tube
undulation vs.\ thick-tube straightness. Caveat: the whole ladder assumes quasi-static mechanics and slow
growth --- if real tip-driven branching (traction, motility) is later needed, the active machinery of the
AVM/SPV (Barton 2017; Bi 2016) is the natural next layer on the 3D geometry.

\paragraph{In sum.} The Turing$\times$vertex \emph{coupling} was never the hard part --- Fig.~4 shows a
clean epithelium carrying any Turing pattern, and the extrusion experiments show morphogen-graded
proliferation localises deformation faithfully ($\mathrm{corr}\!=\!0.42$). The hard part is the
\emph{mechanics substrate}: tubes and branches need shared junction vertices, an enclosed lumen,
quasi-static force balance, and reversible reconnection. The SPV cannot provide them; the parallel Tyssue
vertex model already provides most, and is \emph{three well-defined operators} --- a 3D~T1, an RD-on-mesh
coupling, and the graded-growth port --- from Fig.~5, and \emph{one regime knob} ($\gamma$) beyond that
from Fig.~6.

\end{document}
